\origpage{39}{}
\textbf{Proof.} Let $n = 3p + 2q(3p + c(\Sigma))$. Choose $\theta' > 2q$ so that (8.4) holds, with $\Sigma, p, q, r, \theta$ replaced by $\Sigma, p + n, q, r, \theta'$. Let $\theta = \theta' + n$. We claim that $\theta$ satisfies (9.2). For let $\Gamma, \alpha$ be a portrait in $\Sigma$ of a graph $G$, with depth $\le q$, surrounding a tangle $\mathcal{T}$ of order $\ge \theta$. Suppose that $L$ is a vital $p$-linkage in $G$. Let $\mathcal{A}_1$ be the set of all internal cells $c$ such that $V(\alpha(c)) \setminus \tilde{c}$ contains a terminal of $L$. Since the sets $V(\alpha(c)) \setminus \tilde{c}$ ($c \in \mathcal{A}_1$) are mutually disjoint, by (R7), it follows that $|\mathcal{A}_1| \le p$. By (9.1) and (R9), $\mathcal{A}_1$ contains every cell of $\Gamma$ with $|\tilde{c}| = 3$. By (R6), there are $\le 3|\mathcal{A}_1|$ border cells that are not linked. Let $\mathcal{A}_2 \subseteq C(\Gamma)$ be minimal such that $\mathcal{A}_2$ contains all border cells that are not linked, and contains at least one cell bordering each cuff. Thus, $|\mathcal{A}_2| \le 3p + c(\Sigma)$.

For each internal cell $c \in C(\Gamma) \setminus \mathcal{A}_1$, there is a drawing of $\alpha(c)$ in the closed disc $\bar{c}$ with the vertices in $\tilde{c}$ representing themselves, by (9.1), and this drawing can be chosen so that it meets $bd(\bar{c})$ only in $\tilde{c}$. We may therefore assume, to simplify the notation, that $\alpha(c)$ is such a drawing in $\bar{c}$ for each internal cell $c \in C(\Gamma) \setminus \mathcal{A}_1$. Let $\Gamma'$ be the drawing formed by the vertices in $V(\Gamma)$ and the union of the drawings $\alpha(c)$ over all internal $c \in C(\Gamma) \setminus \mathcal{A}_1$. Then $\Gamma'$ is a drawing in $\Sigma$, and there is a 1-1 correspondence between the spaces of $\Gamma'$ and the border cells of $\Gamma$.

Let $G'$ be the subgraph of $G$ formed by the vertices in $V(\Gamma)$, the graphs $\alpha(u)$ ($u \in V(\Gamma') \cap bd(\Sigma)$) and the graphs $\alpha(c)$ ($c \in C(\Gamma) \setminus (\mathcal{A}_1 \cup \mathcal{A}_2)$). Then $\Gamma' \subseteq G'$. Now $V(\Gamma') \cap bd(\Sigma) = V(\Gamma) \cap bd(\Sigma)$; define $\mu(u) = \alpha(u)$ for each $u \in V(\Gamma') \cap bd(\Sigma)$. Let $F_1^*, \dots, F_r^*$ be the connected components of $bd(\Sigma) \setminus \bigcup(c : c \in \mathcal{A}_2)$. Then $F_1^*, \dots, F_r^*$ are all supports of $\Gamma'$ since $\mathcal{A}_2$ contains at least one cell bordering each cuff; they are mutually disjoint, and $V(\Gamma') \cap bd(\Sigma) \subseteq F_1^* \cup \dots \cup F_r^*$. Moreover, $r = |\mathcal{A}_2| \le 3p + c(\Sigma)$. For each space $s$ of $\Gamma'$ with $s \subseteq F_1^* \cup \dots \cup F_r^*$, let $\mu(s) = \alpha(c)$, where $c$ is the cell of $\Gamma$ with $c \cap bd(\Sigma) = s$.

\hspace*{12pt}(1) $\Gamma', F_1^*, \dots, F_r^*, \mu$ \textit{is a pseudo-presentation of} $G'$.

\textit{Subproof.} Certainly $|V(\Gamma') \cap C| \ge 2$ for each cuff $C$, because $\Gamma$ is a painting and $V(\Gamma) \subseteq V(\Gamma')$. We must verify (P1)-(P4). First let us verify (P1). Now
\begin{align*}
G' &= \Gamma' \cup \bigcup(\alpha(c) : c \in C(\Gamma) \setminus \mathcal{A}_2 \text{ and } c \text{ is a border cell } ) \cup \bigcup(\alpha(u) : u \in V(\Gamma') \cap bd(\Sigma)) \\
&= \Gamma' \cup \mu(F_1^*) \cup \dots \cup \mu(F_r^*).
\end{align*}
By (R3), $E(\mu(s) \cap \mu(s')) = \emptyset$ if $s$ and $s'$ are distinct spaces, since each border cell of $\Gamma$ includes only one space of $\Gamma'$. Since each edge of $\Gamma'$ is an edge of $\alpha(c)$ for some internal cell $c$, it follows from (R3) again that $E(\Gamma' \cap \mu(s)) = \emptyset$ for each space $s \subseteq F_1^* \cup \dots \cup F_r^*$. Finally, let $x_1, x_2$ be units of $\Gamma$ in
